\documentclass[12pt]{article}

% ==================== Fonts (T1 + Latin Modern: full bold/sc typewriter shapes) ====================
\usepackage{lmodern}
\usepackage[T1]{fontenc}
\usepackage[utf8]{inputenc}

% ==================== Core math ====================
\usepackage{amsmath, amssymb, amsthm}
\usepackage{mathtools}

% ==================== Diagrams ====================
\usepackage{tikz}
\usepackage{tikz-cd}
\usetikzlibrary{arrows.meta, positioning, decorations.pathmorphing}

% ==================== Page layout ====================
\usepackage[margin=1in]{geometry}

% Breakable typewriter paths (\path) and ragged table/paragraph columns
\usepackage{url}
\usepackage{array}
\usepackage{placeins}
% Absorb the last few points of over-long lines cleanly (standard arxiv practice)
\setlength{\emergencystretch}{3em}

% ==================== References (hyperref BEFORE cleveref) ====================
\usepackage{graphicx}
\usepackage[colorlinks=true, linkcolor=blue!55!black, citecolor=green!45!black,
            urlcolor=blue!60!black]{hyperref}
\usepackage{cleveref}

% ==================== GrokRxiv sidebar deps ====================
% eso-pic (actively maintained) replaces the deprecated everypage package for
% the page-1 background sidebar.
\usepackage{eso-pic}
\usepackage{textcomp}
\usepackage{xcolor}

% ==================== Theorem environments ====================
\newtheorem{theorem}{Theorem}[section]
\newtheorem{proposition}[theorem]{Proposition}
\newtheorem{lemma}[theorem]{Lemma}
\newtheorem{corollary}[theorem]{Corollary}
\theoremstyle{definition}
\newtheorem{definition}[theorem]{Definition}
\newtheorem{example}[theorem]{Example}
\theoremstyle{remark}
\newtheorem{remark}[theorem]{Remark}
\newtheorem{principle}[theorem]{Principle}

% ==================== Macros ====================
\newcommand{\cat}[1]{\mathcal{#1}}
\newcommand{\Dom}{\mathsf{Dom}}
\newcommand{\Math}{\mathsf{Math}}
\newcommand{\Phys}{\mathsf{Phys}}
\newcommand{\Gpd}{\mathsf{Gpd}}
\newcommand{\Grpd}{\mathbf{Grpd}}
\newcommand{\Set}{\mathbf{Set}}
\newcommand{\Cat}{\mathbf{Cat}}
\newcommand{\Vect}{\mathbf{Vect}}
\newcommand{\Ho}{\mathcal{H}}
\newcommand{\Hphys}{\Ho_{\mathrm{phys}}}
\newcommand{\Hlog}{\Ho_{\mathrm{log}}}
\newcommand{\Hgauge}{\Ho_{\mathrm{gauge}}}
\newcommand{\Herr}{\Ho_{\mathrm{err}}}
\newcommand{\Real}{\mathrm{Real}}
\newcommand{\Obs}{\mathrm{Obs}}
\newcommand{\Repp}{\mathsf{Rep}_{\mathrm{phys}}}
\newcommand{\Aut}{\operatorname{Aut}}
\newcommand{\Stab}{\operatorname{Stab}}
\newcommand{\Orb}{\operatorname{Orb}}
\newcommand{\Hom}{\operatorname{Hom}}
\newcommand{\iHom}{\underline{\operatorname{Hom}}}
\newcommand{\id}{\mathrm{id}}
\newcommand{\op}{\mathrm{op}}
\newcommand{\Spec}{\operatorname{Spec}}
\newcommand{\quot}[2]{[#1/#2]}
\newcommand{\BG}{\mathrm{B}G}
\newcommand{\Pauli}{\mathcal{P}}
\newcommand{\ZZ}{\mathbb{Z}}
\newcommand{\CC}{\mathbb{C}}
\newcommand{\RR}{\mathbb{R}}
\newcommand{\QQ}{\mathbb{Q}}
\newcommand{\FF}{\mathbb{F}}
\newcommand{\antipode}{\mathrm{S}}
\newcommand{\cop}{\Delta}
\newcommand{\MtoP}{Math\textrightarrow{}Physics}
\newcommand{\Ss}{\textsf{S}}
\newcommand{\Hh}{\textsf{H}}
\newcommand{\Pp}{\textsf{P}}
\DeclareMathOperator{\Sheaves}{Sh}
\DeclareMathOperator{\PSh}{PSh}
\DeclareMathOperator{\coker}{coker}
\DeclareMathOperator{\im}{im}

% ==================== GrokRxiv DOI sidebar ====================
\definecolor{grokgray}{RGB}{110,110,110}
\AddToShipoutPictureBG{%
  \ifnum\value{page}=1
    \begin{tikzpicture}[remember picture, overlay]
      \node[
        rotate=90,
        anchor=south,
        font=\Large\sffamily\bfseries\color{grokgray},
        inner sep=0pt
      ] at ([xshift=38pt, yshift=0.52\paperheight]current page.south west)
      {GrokRxiv:2026.07.sheaves-stacks-gauge-quantum-ha\quad
       [\,math.CT\,]\quad
       04 Jul 2026};
    \end{tikzpicture}
  \fi
}

% ==================== Title ====================
\title{\textbf{Sheaves, Stacks, Gauge Redundancy, and\\
Quantum High-Availability}\\[0.4em]
\large Part VI of \emph{A \MtoP{} Representation Library}:\\
The Capstone Module --- Descent, Isotropy, and Fault-Tolerant Encoding}

\author{Matthew Long \\
\textit{The YonedaAI Collaboration} \\
\textit{YonedaAI Research Collective} \\
Chicago, IL \\
\texttt{research@yonedaai.com} $\cdot$ \url{https://yonedaai.com}}

\date{04 July 2026}

\begin{document}
\maketitle

\begin{abstract}
This is Part VI, the capstone module, of the seven-part modular series
\emph{A \MtoP{} Representation Library}. We formalize four intertwined
mathematical domains as physical-representation modules: (i) \textbf{sheaves} as the
theory of compatible local data glued along a Grothendieck topology; (ii)
\textbf{stacks} as fibered categories in groupoids satisfying descent, the objects
that remember not only local data and gluing but also the redundancy symmetries
(automorphisms) of that data; (iii) \textbf{gauge redundancy}, the physical
phenomenon in which many mathematical descriptions represent one physical content,
realized mathematically as a quotient stack $\quot{X}{G}$; and (iv) \textbf{quantum
high-availability}, the fault-tolerant encoding of logical information across
redundant physical degrees of freedom. Our central structural thesis is that these
four are one construction viewed through four lenses: \emph{many local
representatives, one invariant logical content, plus explicit data recording the
redundancy}. We give rigorous definitions and prove four principal results. Theorem
(T1) shows that the quotient stack retains isotropy, $\Aut_{\quot{X}{G}}(x)\simeq
\Stab_G(x)$, and is therefore a strictly more faithful representation object than the
coarse quotient $X/G$ whenever automorphism-counting observables are present. Theorem
(T2), a stackification/descent theorem, discharges the forward reference left open in
Part I: the informally posited ``representation stack'' is literally recovered as the
descent-theoretic stackification of the representation prestack, closing the ladder.
Theorem (T3) proves that surface-code logical operators are the nontrivial homology
classes $H_1(\Sigma;\ZZ_2)$ of the code surface, giving a literal (status \Ss) instance
of ``gauge redundancy $\leftrightarrow$ quantum high-availability'' and reusing Part
IV's homology dictionary verbatim. Theorem (T4) identifies the Connes--Kreimer
antipode with the BPHZ counterterm via Birkhoff decomposition, exhibiting
renormalization as Hopf-algebraic decomposition and reusing Part II's coalgebraic
backbone. We then survey the holographic quantum-error-correction bridge
(Almheiri--Dong--Harlow; Pastawski--Yoshida--Harlow--Preskill; Harlow) as a more
rigorous, operator-algebraic strengthening of the heuristic slogan ``gauge redundancy
is to physics what fault-tolerant encoding is to quantum information.'' Throughout we
preserve the series' three-level epistemic status system (\Ss/\Hh/\Pp) and reproduce
the series limitations, in particular the distinction between gauge redundancy and
physical symmetry. Accompanying Haskell and Lean formalizations encode the
sheaf/stack/stabilizer-code structures.
\end{abstract}

\tableofcontents

% =====================================================================
\section{Introduction}
\label{sec:intro}

\subsection{The modular series and where Part VI sits}

The series \emph{A \MtoP{} Representation Library} advances a single
organizing thesis: many physical quantities are \emph{obtained as realizations}
of structured mathematical information rather than merely \emph{described by} it. The program is deliberately \textbf{modular}: rather than assert a single
grand functor $\Math\to\Phys$, it builds a hierarchy of six standalone, formally
verifiable modules, each introducing one mathematical domain as a
physical-representation faculty and composing on the previous ones, followed by a
seventh synthesis paper. The ladder is
\begin{center}
\resizebox{0.98\linewidth}{!}{$\displaystyle
\underbrace{\text{I}}_{\substack{\text{repr.}\\\text{stack}}}
\to
\underbrace{\text{II}}_{\substack{\text{motives,}\\\text{periods}}}
\to
\underbrace{\text{III}}_{\substack{\text{alg.}\\\text{geometry}}}
\to
\underbrace{\text{IV}}_{\substack{\text{alg.}\\\text{topology}}}
\to
\underbrace{\text{V}}_{\substack{\text{cat.\ theory,}\\\text{HoTT}}}
\to
\underbrace{\text{VI}}_{\substack{\text{sheaves, stacks,}\\\text{gauge, quantum HA}}}
\to
\underbrace{\text{Synth.}}_{\text{recomp.}}
$}
\end{center}
The present paper is \textbf{Part VI}, the capstone before synthesis. Its job is
threefold. First, it introduces the mathematics of \emph{descent}: presheaves,
sheaves, fibered categories, and stacks. Second, it uses that machinery to make
rigorous the two constructions that the whole series has been circling: the
\emph{quotient stack} $\quot{X}{G}$ as the mathematical home of gauge redundancy, and
the \emph{stackification} that converts a prestack into a stack. Third, it exhibits
the principal cross-module bridge in the series: gauge redundancy in physics and
fault-tolerant encoding in quantum information are two readings of one structure,
\emph{many physical representatives, one logical/observable content, with explicit
data recording the redundancy}.

\subsection{The circular closure of the ladder}

A structural fact drives the design. Part I \emph{posits}, by analogy, that
representation entries should glue ``not by equality but by equivalence,'' i.e. that
the assignment of physical representations to mathematical domains behaves like a
\emph{stack}. But Part I does not define a stack; it defers to a forward reference.
Part VI \emph{is} that forward reference resolved: here we define fibered categories,
descent, and stackification, and prove (\Cref{thm:stackification}, T2) that the
representation prestack of Part I admits a universal stackification, satisfying Part
I's informal ``representation stack'' definition exactly. The ladder is therefore
\emph{literally circular}: Part VI supplies the rigorous foundation for Part I's
opening formalism. We flag this as the candidate ``closure theorem'' for the synthesis
paper (Part VII).

\subsection{Contributions}

\begin{enumerate}
\item A self-contained, arxiv-style formalization of sheaves and stacks over a site
(\Cref{sec:sheaves,sec:stacks}), with the physical-representation dictionary of the
series attached to each construction and carrying \Ss/\Hh/\Pp\ status labels.

\item \textbf{T1} (\Cref{prop:stacks-retain}, \Cref{thm:aut-stab}): the quotient stack
retains isotropy, $\Aut_{\quot{X}{G}}(x)\simeq\Stab_G(x)$, and is strictly more
faithful than the coarse quotient for any automorphism-counting observable (orbifold
Euler characteristic, $1/|\Stab|$ path-integral weights).

\item \textbf{T2} (\Cref{thm:stackification}): a descent/stackification theorem that
discharges Part I's forward reference, closing the ladder.

\item \textbf{T3} (\Cref{thm:surface-code}): surface-code logical operators are the
classes of $H_1(\Sigma;\ZZ_2)$; the code space of a genus-$g$ surface code is
$(\CC^2)^{\otimes 2g}$. This is a status-\Ss\ bridge unifying Part IV (homology), Part
VI (stacks/gauge), and quantum information.

\item \textbf{T4} (\Cref{thm:connes-kreimer}): the Connes--Kreimer antipode computes
the BPHZ counterterm and renormalization is the Birkhoff decomposition of the
Feynman-rules character, reusing Part II's Hopf-algebraic decomposition backbone.

\item A survey (\Cref{sec:holographic}) of holographic/operator-algebra quantum error
correction as a rigorous strengthening of the series' central heuristic (status
\Hh$\to$ partially \Ss), with the epistemic status clearly labeled.

\item Companion Haskell (\path{src/sheaves-stacks-gauge-quantum-ha/}) and
Lean (\path{lean/sheaves-stacks-gauge-quantum-ha/}) formalizations.
\end{enumerate}

\subsection{Notation and standing conventions}
\label{subsec:notation}

We work with categories, functors, and (2-)groupoids as developed in Part V. $\Set$ is
the category of sets, $\Grpd$ the (2-)category of groupoids, $\Cat$ that of
(small) categories, $\Vect$ that of vector spaces over a fixed field. For a category
$\cat C$ and object $c$, $\Aut_{\cat C}(c)$ is the automorphism group of $c$. For a
group $G$ acting on a set/space $X$ we write $\Stab_G(x)=\{g\in G: g\cdot x = x\}$ for
the stabilizer, $\Orb_G(x)=\{g\cdot x : g\in G\}$ for the orbit, $X/G$ for the coarse
orbit set, and $\quot{X}{G}$ for the \emph{quotient stack} (\Cref{def:quotient-stack}).
Hilbert spaces are finite-dimensional unless stated otherwise; $\Pauli_n$ is the
$n$-qubit Pauli group. We reserve $\simeq$ for equivalence of groupoids/categories and
$\cong$ for isomorphism of objects.

% =====================================================================
\section{Mathematical framework: the representation stack recalled}
\label{sec:framework}

We briefly recall the load-bearing structures of Parts I--V that Part VI builds on and
discharges; this makes the paper self-contained while fixing notation for the closure
theorem.

\subsection{The realization pipeline and the master formulas}

\begin{definition}[Realization pipeline, after Part I]
\label{def:pipeline}
Let $\Dom$ be a category of \emph{mathematical domains}, $\Phys$ a category of
\emph{physical representations}. For a domain object $M$, the \emph{realization
pipeline} is the composite
\begin{equation}
\Obs_\alpha(M)\;=\;\Obs\bigl(\Real_\alpha(\Phi(M))\bigr),
\label{eq:pipeline}
\end{equation}
where $\Phi(M)$ is the abstract informational content of $M$, $\Real_\alpha$ is a
choice of realization channel (indexed by $\alpha$), and $\Obs$ extracts an
observable. The compressed slogan is $\text{observable}=\Real(\text{abstract
structure})$.
\end{definition}

Part VI instantiates \eqref{eq:pipeline} with $\Real_\alpha$ the \emph{quantum
encoding} $\iota:\Hlog\hookrightarrow\Hphys$ and with $\Obs$ the logical measurement
invariant under the equivalence groupoid of physical variations. The four master
formulas of the series that recur here are
\begin{align}
\Aut_{\quot{X}{G}}(x)&\simeq\Stab_G(x) && \text{(stacks retain isotropy)},
\label{eq:master-aut}\\
\Hphys&\simeq (\Hlog\otimes\Hgauge)\;\oplus\;\Herr &&
\text{(high-availability decomposition)},
\label{eq:master-ha}\\
\cop(\mathcal{A}^{\mathfrak m})&=\textstyle\sum_i
\mathcal{A}^{\mathfrak m}_{i,1}\otimes\mathcal{A}^{\mathfrak m}_{i,2} &&
\text{(coaction / decomposition)},
\label{eq:master-coaction}\\
\delta&:\mathcal{L}\to\textstyle\bigwedge^2\mathcal{L} && \text{(cobracket)}.
\label{eq:master-cobracket}
\end{align}
Equations \eqref{eq:master-aut}--\eqref{eq:master-ha} are proved in
\Cref{sec:stacks,sec:qha}; \eqref{eq:master-coaction}--\eqref{eq:master-cobracket} are
instantiated by the Connes--Kreimer coproduct in \Cref{sec:hopf}.

\subsection{The three-level epistemic status system}

The series' epistemic-honesty device, preserved verbatim, is a three-level status
label attached to every dictionary entry.

\begin{definition}[Status labels \Ss/\Hh/\Pp]
\label{def:status}
An entry of the representation dictionary carries a label:
\begin{itemize}
\item[\Ss] \emph{standard} use in mathematics or mathematical physics (e.g. TQFT as a
functor $\mathrm{Bord}_n\to\Vect$; surface-code logical operators as $H_1(\Sigma;\ZZ_2)$);
\item[\Hh] a strong \emph{heuristic} representation-dictionary entry (e.g. ``gauge
redundancy is fault-tolerant encoding'' as a general slogan);
\item[\Pp] a \emph{speculative} philosophical ontology (e.g. physical reality \emph{as}
motivic-categorical realization).
\end{itemize}
Composite labels \Ss/\Hh\ and \Hh/\Pp\ mark entries standard as pure mathematics but
heuristic/speculative in their physical reading.
\end{definition}

\begin{remark}[Limitation: gauge redundancy is not physical symmetry]
\label{rem:limitation-gauge}
We reproduce the series' governing limitation that is most relevant to Part VI: a
\emph{gauge redundancy} identifies different descriptions of the \emph{same} physical
content, whereas a \emph{physical symmetry} maps one physical state to a
\emph{genuinely different} one. The mathematics of both is a group(oid) action; the
distinction is which action one declares to be ``pure description.'' All status-\Ss\
claims below concern the redundancy reading; the physical-symmetry reading is a
separate modeling choice. This distinction is exactly what the \emph{stack} (versus the
coarse quotient) is built to keep track of.
\end{remark}

% =====================================================================
\section{Sheaves: compatible local data}
\label{sec:sheaves}

\subsection{Sites and presheaves}

\begin{definition}[Grothendieck topology, site]
\label{def:site}
A \emph{Grothendieck topology} $J$ on a small category $\cat C$ assigns to each object
$U\in\cat C$ a collection $J(U)$ of \emph{covering sieves} (equivalently, for the
concrete presentations we use, covering families $\{U_i\to U\}_{i\in I}$), subject to:
\begin{enumerate}
\item (isomorphism) any isomorphism $\{V\xrightarrow{\sim}U\}$ covers $U$;
\item (stability) if $\{U_i\to U\}$ covers and $V\to U$ is any morphism, then the
pullbacks $\{U_i\times_U V\to V\}$ cover $V$;
\item (transitivity) if $\{U_i\to U\}$ covers and each $\{U_{ij}\to U_i\}$ covers, then
$\{U_{ij}\to U\}$ covers.
\end{enumerate}
The pair $(\cat C, J)$ is a \emph{site}. Strictly, the family-level axioms
(isomorphism, stability, transitivity) above present a Grothendieck
\emph{pretopology}; the induced sieve-generated Grothendieck \emph{topology}
(maximality, stability, local character) yields the identical sheaf theory, and we use
the two interchangeably, working throughout with the concrete covering families.
\end{definition}

\begin{example}[The topological site]
\label{ex:top-site}
For a topological space $Y$, let $\mathrm{Op}(Y)$ be the poset of open sets with
inclusions, and let $\{U_i\hookrightarrow U\}$ cover iff $\bigcup_i U_i = U$. This is the
motivating site; ``domains'' $U$ are regions, and a presheaf assigns data to each
region.
\end{example}

\begin{definition}[Presheaf]
\label{def:presheaf}
A \emph{presheaf} on $\cat C$ valued in $\Set$ is a functor
$F:\cat C^{\op}\to\Set$. For $V\to U$ the induced map $F(U)\to F(V)$ is
\emph{restriction}, written $s\mapsto s|_V$. We write $\PSh(\cat C)$ for the category
of presheaves.
\end{definition}

Physically (status \Ss), a presheaf is a \emph{local observable assignment}: to each
region of spacetime, or each patch of a code's interaction hypergraph, it assigns the set of
admissible local configurations, with restriction the operation of forgetting to a
smaller region.

\subsection{The sheaf condition}

\begin{definition}[Sheaf]
\label{def:sheaf}
A presheaf $F$ on a site $(\cat C, J)$ is a \emph{sheaf} if for every covering family
$\{U_i\to U\}\in J(U)$ the diagram
\begin{equation}
\begin{tikzcd}[column sep=large]
F(U) \arrow[r, "e"] &
\displaystyle\prod_i F(U_i)
\arrow[r, shift left=1.1ex, "p"]
\arrow[r, shift right=1.1ex, swap, "q"] &
\displaystyle\prod_{i,j} F(U_i\times_U U_j)
\end{tikzcd}
\label{eq:equalizer}
\end{equation}
is an equalizer, where $e(s)=(s|_{U_i})_i$, $p(s_i)_{ij}=(s_i|_{U_i\times_U U_j})$ and
$q(s_i)_{ij}=(s_j|_{U_i\times_U U_j})$. Concretely: (\emph{separated}) two global
sections agreeing on a cover are equal; (\emph{gluing}) any family of local sections
$s_i\in F(U_i)$ that is \emph{compatible} ($s_i|_{U_{ij}}=s_j|_{U_{ij}}$) descends to a
unique global $s\in F(U)$ with $s|_{U_i}=s_i$.
\end{definition}

\begin{definition}[Stalk, obstruction]
\label{def:stalk}
On the topological site, the \emph{stalk} of $F$ at a point $y$ is the filtered
colimit $F_y=\varinjlim_{U\ni y}F(U)$; it is the germ of local data at the event $y$
(status \Ss). For a sheaf the equalizer \eqref{eq:equalizer} guarantees that compatible
local sections always glue in degree $0$; genuine \emph{obstructions} live in higher
sheaf cohomology $H^{n\ge 1}(U,F)$ --- the failure of a globally defined section to lift
along a surjection of sheaves, or of an abelian-group-valued \v{C}ech $1$-cocycle to be
a coboundary. Such a nonvanishing class represents an anomaly or topological obstruction
(status \Ss/\Hh).
\end{definition}

\begin{proposition}[Sheafification]
\label{prop:sheafification}
The inclusion $\Sheaves(\cat C, J)\hookrightarrow\PSh(\cat C)$ of sheaves into
presheaves admits a left adjoint $(-)^{+\!+}$, the \emph{sheafification}, and the unit
$\eta_F: F\to F^{+\!+}$ is universal among maps from $F$ to sheaves. Sheafification is
obtained by iterating the \emph{plus construction}
$F^+(U)=\varinjlim_{R\in J(U)}\Hom_{\PSh}(R,F)$ twice.
\end{proposition}

\begin{proof}[Proof sketch]
The plus construction $F^+$ is separated for any $F$ and is a sheaf when $F$ is already
separated; hence $F^{+\!+}$ is a sheaf. For any sheaf $G$ and map $f:F\to G$, the
matching-family description of $F^+$ shows $f$ factors uniquely through $\eta$, giving
the adjunction $\Hom_{\Sheaves}(F^{+\!+},G)\cong\Hom_{\PSh}(F,G)$. This is the standard
Grothendieck plus-construction; see Kashiwara--Schapira~\cite{KashiwaraSchapira} and
Vistoli~\cite{Vistoli}. We upgrade it to
groupoid coefficients in \Cref{thm:stackification}.
\end{proof}

The physical content of \Cref{prop:sheafification}: \emph{descent}, reconstructing a
coordinate-independent global field from compatible local measurements, is the
universal solution to the gluing problem. This is the set-valued shadow of the
stack-valued statement we need for gauge fields, to which we now turn.

% =====================================================================
\section{Stacks and gauge redundancy}
\label{sec:stacks}

A sheaf glues compatible local \emph{data}. But a gauge field's local descriptions are
related not by \emph{equality} on overlaps but by \emph{gauge transformations}: on
$U_{ij}$ the two restrictions differ by an element of the gauge group, and these
transition data must satisfy a cocycle condition on triple overlaps. Recording this
requires upgrading from set-valued to groupoid-valued data: a \emph{stack}.

\subsection{Fibered categories and descent}

\begin{definition}[Category fibered in groupoids]
\label{def:cfg}
A functor $\pi:\cat F\to\cat C$ is a \emph{category fibered in groupoids} (CFG) if
(i) for every morphism $f:V\to U$ in $\cat C$ and object $\xi$ over $U$ there is a
\emph{cartesian} arrow $f^*\xi\to\xi$ over $f$, and (ii) every arrow of $\cat F$ is
cartesian. Then each fiber $\cat F(U)=\pi^{-1}(U)$ is a groupoid, and each $f:V\to U$
induces a pullback functor $f^*:\cat F(U)\to\cat F(V)$, functorial up to coherent
natural isomorphism. We write $\cat F(U)$ for the groupoid of objects over $U$.
\end{definition}

Physically, $\cat F(U)$ is the groupoid of local field configurations over the region
$U$, with morphisms the gauge transformations between them (status \Ss).

\begin{definition}[Descent datum, stack]
\label{def:stack}
Let $(\cat C, J)$ be a site and $\pi:\cat F\to\cat C$ a CFG. For a cover
$\{U_i\to U\}$, a \emph{descent datum} is a family $(\xi_i)_i$ with $\xi_i\in\cat F(U_i)$
together with isomorphisms $\varphi_{ij}: \xi_j|_{U_{ij}}\xrightarrow{\sim}
\xi_i|_{U_{ij}}$ satisfying the \emph{cocycle condition}
$\varphi_{ik}=\varphi_{ij}\circ\varphi_{jk}$ on triple overlaps $U_{ijk}$. Descent
data and their morphisms form a groupoid $\mathrm{Desc}(\{U_i\to U\})$. The CFG $\cat F$
is a \emph{stack} if for every cover the natural comparison functor
\begin{equation}
\cat F(U)\;\xrightarrow{\ \simeq\ }\;\mathrm{Desc}(\{U_i\to U\})
\label{eq:descent-equiv}
\end{equation}
is an \emph{equivalence of groupoids} (every descent datum is effective, and this
essentially uniquely). A CFG that satisfies \eqref{eq:descent-equiv} only as a fully
faithful embedding (compatible local objects that glue are unique up to unique iso, but
may fail to exist) is a \emph{prestack}.
\end{definition}

The comparison \eqref{eq:descent-equiv} is precisely the groupoid-valued upgrade of the
sheaf equalizer \eqref{eq:equalizer}: where the sheaf demanded a \emph{unique} glued
section, the stack demands a glued object \emph{unique up to unique isomorphism}, and
in addition demands that the isomorphisms themselves glue.

\begin{center}
\begin{tikzcd}[row sep=large, column sep=huge]
\text{presheaf }F:\cat C^\op\to\Set
  \arrow[r, "\text{coefficients}"]
  \arrow[d, "\text{descent}"']
& \text{CFG }\pi:\cat F\to\cat C
  \arrow[d, "\text{descent}"]\\
\text{sheaf (unique gluing)}
  \arrow[r, "\text{coefficients}"']
& \text{stack (gluing up to unique iso)}
\end{tikzcd}
\end{center}

\subsection{The quotient stack and isotropy}

\begin{definition}[Action groupoid and quotient stack]
\label{def:quotient-stack}
Let a group $G$ act on a space $X$. The \emph{action groupoid} $X/\!\!/G$ has objects
the points of $X$ and morphisms $x\xrightarrow{g}g\cdot x$ for $g\in G$, with
composition given by group multiplication. The \emph{quotient stack} $\quot{X}{G}$ is
the stackification of the prestack $U\mapsto \{G\text{-torsors }P\to U \text{ with
}G\text{-map } P\to X\}$; over a point it is the groupoid $X/\!\!/G$. The \emph{coarse
quotient} $X/G$ is the set of orbits $\pi_0(X/\!\!/G)$.
\end{definition}

The two quotients are related by the coarse-space map $c:\quot{X}{G}\to X/G$, which is
initial among maps from $\quot{X}{G}$ to sets/algebraic spaces. It is a bijection on
isomorphism classes of points but collapses all automorphisms.

\begin{proposition}[Stacks retain gauge information --- T1, qualitative]
\label{prop:stacks-retain}
If a physical configuration space is represented by descriptions $X$ modulo gauge
transformations $G$, and if stabilizer groups influence observables (anomalies,
quantization, path-integral weights), then $\quot{X}{G}$ is a strictly more faithful
representation object than $X/G$.
\end{proposition}

\begin{proof}
The coarse quotient records orbits but forgets automorphism groups of representatives:
$c(x)=c(x')$ iff $\Orb_G(x)=\Orb_G(x')$, and $c$ carries no further data. The quotient
stack records objects, isomorphisms (gauge transformations), and self-isomorphisms
(residual/isotropy symmetry). By \Cref{thm:aut-stab} below, the self-isomorphisms at
$x$ are exactly $\Stab_G(x)$. Any observable that is a function of the isotropy group
--- e.g. an orbifold Euler characteristic $\chi_{\mathrm{orb}}=\sum_{[x]}
1/|\Stab_G(x)|$ or a path-integral measure weighting each orbit by $1/|\Stab_G(x)|$ ---
is therefore computable from $\quot{X}{G}$ but not from $X/G$, which has discarded the
$|\Stab_G(x)|$ factors. Hence $\quot{X}{G}\to X/G$ is a genuine loss of information
whenever some $\Stab_G(x)$ is nontrivial and enters an observable.
\end{proof}

\begin{theorem}[Isotropy of the quotient stack --- T1, quantitative]
\label{thm:aut-stab}
For a $G$-action on $X$ and any point $x\in X$, the automorphism group of $x$ regarded
as an object of $\quot{X}{G}$ is naturally isomorphic to the stabilizer:
\begin{equation}
\Aut_{\quot{X}{G}}(x)\;\cong\;\Stab_G(x).
\label{eq:aut-stab}
\end{equation}
Moreover the \emph{residual gerbe} of $\quot{X}{G}$ at $[x]$ --- the restriction of the
stack to the orbit of $x$ --- is the classifying stack
$\mathrm{B}\Stab_G(x)=\quot{\mathrm{pt}}{\Stab_G(x)}$. (The full \'etale-local model of
$\quot{X}{G}$ near $[x]$ is the slice quotient $\quot{S_x}{\Stab_G(x)}$ for a
$\Stab_G(x)$-invariant normal slice $S_x$ to the orbit, reducing to
$\mathrm{B}\Stab_G(x)$ exactly when the orbit is open.)
\end{theorem}

\begin{proof}
Work in the action groupoid $X/\!\!/G$, which computes the fiber of $\quot{X}{G}$ over a
point (\Cref{def:quotient-stack}). By definition a morphism $x\to x'$ in $X/\!\!/G$ is an
element $g\in G$ with $g\cdot x = x'$, and composition is group multiplication. An
\emph{automorphism} of $x$ is a morphism $x\to x$, i.e. an element $g\in G$ with
$g\cdot x = x$; this is exactly the condition $g\in\Stab_G(x)$. The bijection
$\Aut_{X/\!\!/G}(x)=\{g\in G: g\cdot x=x\}=\Stab_G(x)$ is a group isomorphism because
composition of automorphisms is group multiplication, matching the group law on
$\Stab_G(x)$. This proves \eqref{eq:aut-stab}.

For the residual gerbe: restrict the action groupoid to the single orbit of $x$. The
orbit $\Orb_G(x)\cong G/\Stab_G(x)$ contributes no automorphisms (it is a free
direction), while the residual symmetry is generated by $\Stab_G(x)$. The restricted
groupoid is the transitive groupoid on $G/\Stab_G(x)$, which is equivalent as a groupoid
to the one-object groupoid with automorphism group $\Stab_G(x)$, i.e. to
$\mathrm{B}\Stab_G(x)$. Hence the restriction of $\quot{X}{G}$ to the orbit --- its
residual gerbe at $[x]$ --- is $\mathrm{B}\Stab_G(x)$; the full \'etale-local model
additionally records the normal slice as $\quot{S_x}{\Stab_G(x)}$.
\end{proof}

\begin{corollary}[Faithfulness gap is exactly the isotropy]
\label{cor:faithfulness}
The fibers of $c:\quot{X}{G}\to X/G$ over $[x]$ are the classifying stacks
$\mathrm{B}\Stab_G(x)$; equivalently, the information lost by passing from the stack to
the coarse quotient is precisely the groupoid cohomology
$H^\bullet(\mathrm{B}\Stab_G(x))=H^\bullet_{\mathrm{grp}}(\Stab_G(x))$ of the isotropy
groups. In particular $X/G$ is faithful iff the action is free.
\end{corollary}

\begin{proof}
Immediate from \Cref{thm:aut-stab}: over the orbit $[x]$ the coarse map forgets exactly
the one-object groupoid $\mathrm{B}\Stab_G(x)$, whose invariants are the group
cohomology of $\Stab_G(x)$. Freeness means all $\Stab_G(x)$ are trivial, so $c$ is an
equivalence.
\end{proof}

\begin{example}[Electromagnetism as gauge redundancy]
\label{ex:em}
Let $X$ be the space of $U(1)$ connections (gauge potentials $A$) on a manifold and
$G$ the group of gauge transformations acting by $A\mapsto A+d\lambda$. The physical
content is the curvature $F=dA$, and $d^2=0$ gives $F$ gauge invariance. The coarse
quotient $X/G$ records the physical field strengths. The stack $\quot{X}{G}$ in
addition records the automorphisms: for a generic connection the stabilizer is the
constant gauge transformations $U(1)$, and this residual $U(1)$ isotropy is exactly the
global symmetry that survives quantization and contributes the $1/\mathrm{vol}(G)$
Faddeev--Popov factors to the path integral. This is the canonical status-\Ss\
illustration of \Cref{prop:stacks-retain}.
\end{example}

\begin{remark}[Classifying stack $\mathrm{B}G$]
\label{rem:BG}
The special case $X=\mathrm{pt}$ gives $\quot{\mathrm{pt}}{G}=\mathrm{B}G$, the
classifying stack, whose groupoid of points over any base is the groupoid of
$G$-torsors. It classifies $G$-bundles (gauge fields), so the dictionary entry
``classifying stack $\mathrm{B}G$ $\leftrightarrow$ universal $G$-bundle object / gauge
field classifier'' (status \Ss) is literally the statement that
$\Hom(-,\mathrm{B}G)$ is the groupoid of gauge fields.
\end{remark}

% =====================================================================
\section{Stackification and the closure of the ladder}
\label{sec:closure}

We now discharge Part I's forward reference. Part I posited a \emph{representation
prestack} $\Repp:\Dom^{\op}\to\Grpd$ assigning to each mathematical domain the
groupoid of its physical representations, and asserted, by analogy, that this
should be a \emph{stack} (local representations glue up to gauge equivalence). Part I
did not prove this. Here we do, via stackification.

\begin{theorem}[Stackification of the representation prestack --- T2]
\label{thm:stackification}
Let $(\Dom, J)$ be a site and $\Repp:\Dom^{\op}\to\Grpd$ the representation prestack of
Part I (a CFG over $\Dom$). Then:
\begin{enumerate}
\item There exists a stack $\Repp^{J}$ on $(\Dom, J)$ and a morphism of CFGs
$\eta:\Repp\to\Repp^{J}$ such that for every stack $\cat G$ on $(\Dom,J)$, precomposition
\[
\eta^*:\ \iHom_{\mathrm{Stack}}(\Repp^{J},\cat G)\ \xrightarrow{\ \simeq\ }\
\iHom_{\mathrm{CFG}}(\Repp,\cat G)
\]
is an equivalence of groupoids. That is, $\Repp^{J}$ is the universal stack under
$\Repp$ (\emph{stackification}).
\item $\Repp$ is already a stack iff $\eta$ is an equivalence.
\item Consequently, Part I's ``representation stack'' exists and is well-defined: it is
$\Repp^{J}$. When the local representation data of Parts II--V (Betti/de Rham/Hodge
realizations, cohomology, TQFT functors, HoTT identity data) are compatible on
overlaps, they assemble into a global object of $\Repp^{J}$, exactly satisfying Part I's
informal definition ``entries glue not by equality but by gauge equivalence.''
\end{enumerate}
\end{theorem}

\begin{proof}
We adapt the plus construction of \Cref{prop:sheafification} from set-valued to
groupoid-valued (equivalently, $(2,1)$-truncated) coefficients. Define, for
$U\in\Dom$,
\[
\Repp^{+}(U)\;:=\;\varinjlim_{\{U_i\to U\}\in J(U)}\ \mathrm{Desc}\bigl(\{U_i\to
U\},\Repp\bigr),
\]
the (pseudo-)colimit over covers of the descent groupoids of \Cref{def:stack}, with
transition functors the refinement functors of covers. This colimit is filtered: any two
covers of $U$ admit a common refinement (their pairwise pullbacks $\{U_i\times_U
V_j\to U\}$), so the indexing category of covers is directed and the (pseudo-)colimit is
well-behaved. Two applications, $\Repp^{J}:=
\Repp^{+\!+}$, produce a stack:
\emph{(a)} $\Repp^{+}$ is a prestack (fully faithful descent) for any CFG $\Repp$,
because morphisms in a descent groupoid are already determined locally and the cocycle
condition makes them glue; \emph{(b)} if $\Repp$ is a prestack then $\Repp^{+}$ is a
stack, because effectivity of descent is achieved after one further plus step, exactly
as in the $1$-categorical case but now tracking the coherence isomorphisms
$\varphi_{ij}$ and their cocycle condition. Universality: given a stack $\cat G$ and a
map $f:\Repp\to\cat G$, descent in $\cat G$ (\Cref{def:stack}) means every descent
datum in $\cat G$ is effective, so $f$ extends over each plus step uniquely up to
unique coherent isomorphism; this yields the essentially unique factorization through
$\eta$, i.e. the claimed equivalence of $\Hom$-groupoids. Part (2) is immediate: if
$\eta$ is an equivalence then $\Repp\simeq\Repp^{J}$ is a stack, and conversely a stack
is $J$-local so the plus construction is idempotent on it. Part (3) is the application:
the compatibility of the local realization data supplied by Parts II--V is exactly a
descent datum for $\Repp$, hence an object of $\Repp^{J}(U)$; this is the content of Part
I's informal gluing axiom, now a theorem. This is the general stackification theorem for
fibered categories over a site (Giraud~\cite{Giraud}; Vistoli~\cite{Vistoli}),
specialized to $\Repp$.
\end{proof}

\begin{center}
\begin{tikzcd}[row sep=large, column sep=huge]
\Repp \arrow[r, "\eta"] \arrow[dr, swap, "\forall f"'] &
\Repp^{J} \arrow[d, dashed, "\exists! \bar f \text{ (up to iso)}"] \\
& \cat G \ (\text{any stack})
\end{tikzcd}
\end{center}

\begin{corollary}[Closure of the ladder]
\label{cor:closure}
The six modules I--VI are not merely an expository decomposition: \Cref{thm:stackification}
shows that Part I's opening definition (the representation stack) is a theorem provable
from Part VI's descent machinery applied to the concrete data of Parts II--V. The ladder
$\text{I}\to\cdots\to\text{VI}$ therefore closes back onto Part I. We nominate this as the
headline ``closure theorem'' for the synthesis paper (Part VII).
\end{corollary}

\begin{remark}[Status]
\label{rem:t2-status}
\Cref{thm:stackification}(1)--(2) is status \Ss\ (the stackification theorem is standard;
Giraud~\cite{Giraud}, Vistoli~\cite{Vistoli}). Part (3), the identification of the
specific descent datum with
Part I's ``representation stack,'' is status \Ss/\Hh: the mathematics is standard but the
claim that the physical realization data of Parts II--V \emph{is} a descent datum is the
series' modeling hypothesis, and we flag it as such.
\end{remark}

% =====================================================================
\section{Quantum high-availability encodings}
\label{sec:qha}

We now cross from geometry to quantum information. The same slogan (many physical
representatives, one logical content, plus explicit redundancy data) becomes the
theory of quantum error-correcting codes. We make the analogy precise and then, in
\Cref{sec:surface}, exhibit a case where it is a literal (status \Ss) identity.

\subsection{Codes as encodings with an equivalence groupoid}

\begin{definition}[Quantum high-availability representation]
\label{def:qha}
A \emph{quantum high-availability (QHA) representation} is an isometric encoding
$\iota:\Hlog\hookrightarrow\Hphys$ together with a groupoid $\mathcal E$ of
\emph{physical variations} (channels acting on $\Hphys$) such that:
\begin{enumerate}
\item every logical observable $\overline{O}$ (an operator on $\iota(\Hlog)$) is
\emph{invariant} under $\mathcal E$ in the sense that its induced action on the code
space is unchanged by any variation in $\mathcal E$; and
\item detectable/correctable errors move the system within controlled equivalence or
syndrome sectors, from which a recovery channel returns it to $\iota(\Hlog)$.
\end{enumerate}
The redundancy $\dim\Hphys/\dim\Hlog$ is the ``availability'': logical information
survives loss of any physical subsystem in the correctable set.
\end{definition}

This is the realization pipeline \eqref{eq:pipeline} with $\Real=\iota$ and $\Obs$ the
logical measurement; $\mathcal E$ plays the role of the gauge groupoid of
\Cref{sec:stacks}. The subsystem structure is the high-availability decomposition
\eqref{eq:master-ha}:
\begin{equation}
\Hphys\ \simeq\ (\Hlog\otimes\Hgauge)\ \oplus\ \Herr,
\label{eq:subsystem}
\end{equation}
where $\Hgauge$ carries \emph{gauge qubits} (redundant degrees of freedom that may be
changed freely without affecting logical content) and $\Herr$ is the error sector
detected by syndrome measurement. We stress the distinction between the two ways
redundancy is realized. In a genuine \emph{subsystem} (Bacon--Shor--type) or
operator-algebra code the gauge factor is nontrivial, $\dim\Hgauge>1$, and gauge
operators act nontrivially within $\Hlog\otimes\Hgauge$. In a \emph{pure stabilizer}
code the gauge factor degenerates, $\dim\Hgauge=1$: the stabilizers act as the identity
on the code space and the redundancy is carried entirely by the orthogonal syndrome
sectors $\Herr$, i.e. by the many ways the full kinematic space projects onto
$\iota(\Hlog)$. The parallel with $\quot{X}{G}$ is then this: $\Hlog=$ logical/observable
content (analogue of $X/G$), the redundancy the stack remembers ($\Stab$/isotropy)
appears as $\Hgauge$ for subsystem codes and as the $\Herr$ syndrome grading for pure
stabilizer codes, and in both cases the encoding $\iota$ is the faithful (stack-level)
representative that the coarse logical content alone forgets.

\subsection{Stabilizer and operator-algebra codes}

\begin{definition}[Stabilizer code]
\label{def:stabilizer}
Let $\Pauli_n$ be the $n$-qubit Pauli group. A \emph{stabilizer code}~\cite{Gottesman} is
specified by an abelian subgroup $S\subset\Pauli_n$ containing no nontrivial phases,
$S\cap\{\pm I,\pm iI\}=\{I\}$ (in particular $-I\notin S$, so the code space is
nonzero); the \emph{code space} is the simultaneous $+1$ eigenspace
\begin{equation}
\Hlog\ =\ \{\,|\psi\rangle\in(\CC^2)^{\otimes n}\ :\ s|\psi\rangle=|\psi\rangle\ \ \forall
s\in S\,\}.
\label{eq:code-space}
\end{equation}
If $S$ has $n-k$ independent generators then $\dim\Hlog=2^{k}$ (it encodes $k$ logical
qubits). The \emph{logical operators} are elements of the normalizer
$N(S)/S\subset\Pauli_n/S$: Pauli operators commuting with all of $S$ but not in $S$.
\end{definition}

\begin{lemma}[Knill--Laflamme conditions]
\label{lem:knill-laflamme}
Let $P$ be the projector onto a code space $\Hlog\subset\Hphys$ and $\{E_a\}$ a set of
error operators. Then $\{E_a\}$ is correctable iff
\begin{equation}
P E_a^\dagger E_b P\ =\ c_{ab}\,P
\label{eq:kl}
\end{equation}
for a Hermitian matrix $(c_{ab})$. In words: the environment learns nothing about the
logical state from the errors (no logical information leaks into which-error data), so a
recovery channel exists.
\end{lemma}

\begin{proof}[Proof sketch]
Necessity: if a recovery channel $\mathcal R$ satisfies $\mathcal
R(E_a\rho E_b^\dagger)\propto\rho$ for all code states $\rho=P\rho P$, tracing against
$P$ forces \eqref{eq:kl}. Sufficiency: diagonalize $(c_{ab})=u\,d\,u^\dagger$; the
combinations $F_\mu=\sum_a u_{a\mu}E_a$ satisfy $PF_\mu^\dagger F_\nu P=d_\mu
\delta_{\mu\nu}P$, so the $F_\mu$ map $\Hlog$ to mutually orthogonal subspaces
isometrically (up to normalization); measuring which subspace (the \emph{syndrome}) and
inverting the corresponding isometry recovers the logical state. This is the standard
Knill--Laflamme theorem~\cite{KnillLaflamme}.
\end{proof}

\begin{proposition}[Stabilizer syndromes form a sheaf; correction is the descent --- T3, part i]
\label{prop:syndrome-presheaf}
For a stabilizer code with interaction hypergraph $\Gamma$ (vertices = qubits,
hyperedges = stabilizer supports, since a weight-$w$ generator such as a surface-code
star or plaquette touches $w>2$ qubits; equivalently the bipartite Tanner graph with
qubit- and check-nodes), the assignment
\[
U\ \longmapsto\ \{\text{syndrome outcomes } \sigma:S_U\to\{\pm1\}\text{ consistent with
} S\},\qquad U\subseteq\Gamma,
\]
where $S_U$ are the stabilizer generators supported in $U$, is a sheaf on the poset
of sub-hypergraphs, with restriction given by forgetting outcomes outside $U$. The
substantive error-correction content is a \emph{second} descent, one level up: the map
\[
(\text{global syndrome } \sigma:S\to\{\pm1\})\ \longmapsto\ (\text{correctable global
error class in }\Pauli_n/S),
\]
is well-defined precisely when the Knill--Laflamme conditions (\Cref{lem:knill-laflamme})
hold on the correctable set.
\end{proposition}

\begin{proof}
Restriction of syndrome outcomes to a sub-hypergraph is functorial and strictly
compatible with further restriction, so the assignment is a presheaf
(\Cref{def:presheaf}); moreover a family of local syndrome functions that agree on
overlaps glues uniquely to a global syndrome function, since a syndrome is just an
assignment of $\pm1$ to each generator --- so syndromes trivially form a \emph{sheaf}.
This gluing is automatic and carries no error-correction content by itself. The
non-trivial statement is the well-definedness of the recovery map: a globally consistent
syndrome $\sigma$ determines a \emph{unique} correctable error class in $\Pauli_n/S$
(equivalently a unique recovery up to a stabilizer) exactly when the Knill--Laflamme
matrix \eqref{eq:kl} is nondegenerate on the correctable set. Thus the sheaf of syndromes
is the easy layer, and correctability is the genuine descent from global syndrome data to
a unique global logical correction.
\end{proof}

\begin{remark}[Operator-algebra generalization]
\label{rem:oaqec}
The subsystem decomposition \eqref{eq:subsystem} is the framework of
\emph{operator-algebra quantum error correction} (OAQEC): the correctable code is a von
Neumann subalgebra $\mathcal A\subset\mathcal B(\Hphys)$, and logical operators are
$\mathcal A$, gauge operators its commutant restricted to the code, errors the
complement. When $\mathcal A$ has nontrivial center one recovers gauge theories with a
center (edge modes); this is the setting of \Cref{sec:holographic}.
\end{remark}

% =====================================================================
\section{Surface codes: gauge redundancy is literally homology}
\label{sec:surface}

The analogy of \Cref{sec:qha} becomes a theorem for topological codes. Here the
logical content is a \emph{homology group} of the code surface: Part IV's dictionary
entry ``homology $H_n(X)$ $\leftrightarrow$ conserved extended structures'' realized
literally, not analogically. This is the strongest status-\Ss\ bridge in the series.

\begin{definition}[Surface code]
\label{def:surface-code}
The surface/toric code~\cite{Kitaev,DKLP} is defined as follows. Let $\Sigma$ be a closed
\emph{connected} orientable surface of genus $g$ with a cellulation (a lattice embedded in
$\Sigma$).
Place a qubit on each edge. For each vertex $v$ define the
\emph{star operator} $A_v=\prod_{e\ni v}X_e$ (product over edges incident to $v$); for
each face $f$ define the \emph{plaquette operator} $B_f=\prod_{e\in\partial f}Z_e$
(product over edges in the boundary of $f$). The stabilizer group is
$S=\langle A_v, B_f\rangle$. These all commute (each $A_v$ and $B_f$ share an even number
of edges), so \eqref{eq:code-space} defines the code space.
\end{definition}

\begin{theorem}[Surface-code logical operators are the (co)homology $H_1, H^1(\Sigma;\ZZ_2)$ --- T3]
\label{thm:surface-code}
For the surface code on a genus-$g$ surface $\Sigma$:
\begin{enumerate}
\item the code space has dimension $2^{2g}$, i.e. $\Hlog\cong(\CC^2)^{\otimes 2g}$ and
the code encodes $k=2g$ logical qubits;
\item the $Z$-type and $X$-type logical operators (modulo stabilizers) are respectively
the nontrivial classes of $H_1(\Sigma;\ZZ_2)$ and the Poincar\'e-dual classes of
$H^1(\Sigma;\ZZ_2)$, each isomorphic to $\ZZ_2^{2g}$, so that the full logical Pauli
group modulo stabilizers is the $4g$-dimensional symplectic $\ZZ_2$-space
\[
N(S)/S\ \cong\ H_1(\Sigma;\ZZ_2)\ \oplus\ H^1(\Sigma;\ZZ_2)\ \cong\ \ZZ_2^{\,4g};
\]
each homological summand contributes the $2g$ generators of one Pauli type for the $2g$
logical qubits, and the intersection pairing supplies the symplectic (anti)commutation
form;
\item error strings that are boundaries $c=\partial b$ act trivially on the code space
(they are ``pure gauge''); only homologically nontrivial (noncontractible) strings
implement logical operations.
\end{enumerate}
\end{theorem}

\begin{proof}
\emph{(1) Dimension count.} There are $|E|$ qubits, and $|V|+|F|$ stabilizer generators
$A_v, B_f$. They are not all independent: $\prod_v A_v = I$ and $\prod_f B_f = I$ (each
edge appears in exactly two stars and two plaquettes over $\ZZ_2$), and these are the
\emph{only} relations --- this holds precisely because $\Sigma$ is a \emph{closed}
surface (no boundary), so every edge is shared by exactly two faces and two vertices.
Hence the number of independent generators is
$(|V|-1)+(|F|-1)=|V|+|F|-2$. The number of encoded qubits is
\[
k\ =\ |E|-(|V|+|F|-2)\ =\ 2-(|V|-|E|+|F|)\ =\ 2-\chi(\Sigma)\ =\ 2-(2-2g)\ =\ 2g,
\]
using the Euler characteristic $\chi(\Sigma)=|V|-|E|+|F|=2-2g$. Thus $\dim\Hlog=2^{2g}$.

\emph{(2) Homological identification.} Work over $\ZZ_2$. A $Z$-type Pauli
$\prod_{e\in c}Z_e$ is indexed by a $1$-chain $c\in C_1(\Sigma;\ZZ_2)$. It commutes with
every star $A_v$ iff $c$ is a \emph{cycle} ($\partial c=0$), since the number of edges of
$c$ at each vertex must be even. It lies in the stabilizer $S$ (is a product of
plaquettes) iff $c$ is a \emph{boundary} ($c=\partial b$ for a $2$-chain $b$). Hence
the $Z$-logical operators modulo stabilizers are
\[
\ker\partial_1/\im\partial_2\ =\ H_1(\Sigma;\ZZ_2)\ \cong\ \ZZ_2^{2g}.
\]
The same argument on the dual cellulation gives the $X$-logical operators as
$H_1(\Sigma^\vee;\ZZ_2)\cong H^1(\Sigma;\ZZ_2)\cong\ZZ_2^{2g}$, Poincar\'e-dual to the
$Z$-cycles. The full logical Pauli group modulo stabilizers is therefore the direct sum
$N(S)/S\cong H_1\oplus H^1\cong\ZZ_2^{4g}$; the $\ZZ_2$-valued intersection pairing on
$H_1\times H^1$ supplies the nondegenerate symplectic form implementing the canonical
(anti)commutation of the $2g$ logical qubit pairs (each qubit an $X_i,Z_i$ pair drawn one
from each homological summand).

\emph{(3) Boundaries are pure gauge.} If $c=\partial b$ then $\prod_{e\in c}Z_e=\prod_{f\in
b}B_f\in S$, so it acts as the identity on the code space \eqref{eq:code-space}. Thus a
contractible error loop is undetectable and harmless, while only noncontractible loops
--- generators of $H_1$ --- change the logical state. This is precisely the AT-library
entry ``boundary $c=\partial b$ $\leftrightarrow$ pure gauge or physically null
sector.''
\end{proof}

\begin{center}
\begin{tikzcd}[row sep=large, column sep=large]
C_2(\Sigma;\ZZ_2) \arrow[r, "\partial_2"] &
C_1(\Sigma;\ZZ_2) \arrow[r, "\partial_1"] &
C_0(\Sigma;\ZZ_2)\\
\text{plaquettes } B_f \arrow[u, phantom, "\rotatebox{90}{$\in$}"] &
\text{$Z$-strings} \arrow[u, phantom, "\rotatebox{90}{$\in$}"] &
\text{stars } A_v \arrow[u, phantom, "\rotatebox{90}{$\in$}"]
\end{tikzcd}
\end{center}

\begin{corollary}[Gauge redundancy $\leftrightarrow$ quantum high availability, literal case]
\label{cor:literal-bridge}
The surface code realizes \eqref{eq:subsystem} and \eqref{eq:aut-stab} simultaneously.
The logical content is the homology $\Hlog\cong(\CC^2)^{\otimes 2g}$ with logical
operators graded by $H_1(\Sigma;\ZZ_2)\cong\ZZ_2^{2g}$. Because the surface code is a
\emph{pure} stabilizer code, the gauge factor is trivial ($\dim\Hgauge=1$): the
contractible-loop (boundary) operators lie \emph{in} the stabilizer $S$ and act as the
identity on the code space, so they are the ``pure gauge'' null sector rather than
operators on a nontrivial tensor factor. The genuine redundancy is therefore carried by
the syndrome grading $\Herr$: the violated stars/plaquettes $A_v, B_f$ label the
orthogonal error sectors through which the full $2^{|E|}$-dimensional kinematic space
projects onto the $2^{2g}$-dimensional code space. (A Bacon--Shor--type subsystem code on
$\Sigma$ would instead populate a nontrivial $\Hgauge$; see the remark on OAQEC below.)
Hence ``gauge redundancy is to physics what fault-tolerant encoding is to quantum
information'' is, \emph{for surface codes}, a theorem (status \Ss), not a slogan: the
gauge-null (boundary) sector and the logical sector are the two homological strata
$\mathrm{im}\,\partial_2$ and $H_1(\Sigma;\ZZ_2)$ of one and the same chain complex.
\end{corollary}

\begin{remark}[Cross-module reuse]
\label{rem:crossmodule-t3}
\Cref{thm:surface-code} reuses Part IV's homology dictionary verbatim
($H_n(X)\leftrightarrow$ conserved extended structures) and instantiates Part I's
Equivalence axiom (gauge-equivalent = homologous). It is the module's central case
study and the reason Part VI has the most cross-references of any module.
\end{remark}

% =====================================================================
\section{Hopf-algebraic decomposition: renormalization as antipode}
\label{sec:hopf}

The second reused backbone is coalgebraic. Part II introduced the Goncharov~\cite{Goncharov} Lie
coalgebra and the coaction \eqref{eq:master-coaction} as the ``anatomy'' of amplitudes.
Part VI uses the \emph{same} formal operation (a graded connected Hopf-algebra
coproduct) on a different combinatorial datatype (Feynman graphs / rooted trees),
where the antipode computes renormalization counterterms.

\begin{definition}[Hopf algebra of rooted trees]
\label{def:ck}
Let $H$ be the free commutative algebra over $\QQ$ on rooted trees (forests as products),
graded by the number of vertices (loop number). The \emph{coproduct} is the sum over
admissible cuts:
\begin{equation}
\cop(\Gamma)\ =\ \Gamma\otimes 1\ +\ 1\otimes\Gamma\ +\ \sum_{\text{cuts } c}
P^c(\Gamma)\otimes R^c(\Gamma),
\label{eq:ck-coproduct}
\end{equation}
where $R^c$ is the component containing the root and $P^c$ the pruned forest. $H$ is a
graded connected bialgebra; connectedness ($H_0=\QQ$) forces a unique \emph{antipode}
$\antipode$ making $H$ a Hopf algebra, given recursively by
\begin{equation}
\antipode(\Gamma)\ =\ -\Gamma\ -\ \sum_{\text{cuts } c}\antipode\bigl(P^c(\Gamma)\bigr)\,
R^c(\Gamma).
\label{eq:ck-antipode}
\end{equation}
\end{definition}

\begin{definition}[Feynman-rules character and Birkhoff decomposition]
\label{def:character}
A regularized \emph{Feynman-rules character} is an algebra map
$\varphi:H\to\CC(\!(z)\!)$ (formal Laurent series in the dimensional-regularization parameter
$z$), an element of the group $G=\Hom_{\mathrm{alg}}(H,\mathcal A)$ of $\mathcal
A$-valued characters under the convolution product $\varphi\star\psi=m_{\mathcal
A}\circ(\varphi\otimes\psi)\circ\cop$. The minimal-subtraction splitting $\mathcal
A=\mathcal A_-\oplus\mathcal A_+$ (pole part vs. holomorphic part) is a Rota--Baxter
structure of weight $-1$.
\end{definition}

\begin{theorem}[Antipode = counterterm; renormalization = Birkhoff decomposition --- T4]
\label{thm:connes-kreimer}
Let $\varphi:H\to\mathcal A$ be a Feynman-rules character valued in a commutative
Rota--Baxter algebra $\mathcal A=\mathcal A_-\oplus\mathcal A_+$ with projection $R$ onto
$\mathcal A_-$. Then $\varphi$ admits a unique \emph{Birkhoff decomposition}
$\varphi=\varphi_-^{\star-1}\star\varphi_+$ in $G$, with the \emph{counterterm}
$\varphi_-$ and \emph{renormalized character} $\varphi_+$ given recursively by Bogoliubov's
formulas
\begin{align}
\varphi_-(\Gamma)\ &=\ -R\Bigl(\varphi(\Gamma)+\sum \varphi_-(P^c\Gamma)\,\varphi(R^c\Gamma)\Bigr),
\label{eq:bphz-minus}\\
\varphi_+(\Gamma)\ &=\ (1-R)\Bigl(\varphi(\Gamma)+\sum \varphi_-(P^c\Gamma)\,\varphi(R^c\Gamma)\Bigr).
\label{eq:bphz-plus}
\end{align}
The counterterm $\varphi_-=\antipode_R^{\varphi}$ is the twisted (Rota--Baxter) antipode
--- an operator on the character group $\Hom(H,\mathcal A)$, not a precomposition with a
map on $H$, since $R$ acts on the target algebra $\mathcal A$ --- and $\varphi_+$ is the
BPHZ-renormalized amplitude. Thus renormalization is exactly the Hopf-algebraic
decomposition \eqref{eq:master-coaction} applied to the divergence structure.
\end{theorem}

\begin{proof}[Proof sketch]
Since $H$ is graded connected, the convolution group $G$ has a well-defined exponential/
logarithm and every character has a $\star$-inverse computed by the antipode
\eqref{eq:ck-antipode}. Writing $\varphi_-=e\circ\epsilon - R\circ(\varphi_-\star(\varphi-e\circ\epsilon))$
and solving degree by degree gives \eqref{eq:bphz-minus}; the Rota--Baxter
identity~\cite{RotaBaxter}
\[
R(a)R(b)+R(ab)=R\bigl(R(a)b+aR(b)\bigr)
\]
is precisely what makes $\varphi_-$ multiplicative (a character), so the decomposition
lands in the group $G$. The holomorphic part $\varphi_+=\varphi_-\star\varphi$ is then
finite at $z=0$ and equals the forest formula of BPHZ. Identifying the recursion with
the antipode recursion \eqref{eq:ck-antipode} twisted by $R$ gives
$\varphi_-=\antipode_R^{\varphi}$, the twisted antipode on $\Hom(H,\mathcal A)$. This is
the Connes--Kreimer Riemann--Hilbert theorem~\cite{ConnesKreimer}.
\end{proof}

\begin{corollary}[Shared Hopf backbone of Parts II and VI]
\label{cor:shared-hopf}
Both the Goncharov coaction (Part II) and the Connes--Kreimer coproduct (Part VI) are
coproducts of graded connected Hopf algebras; they are the same formal operation
\eqref{eq:master-coaction} applied to two combinatorial datatypes (polylog symbols vs.
rooted trees), with two physical readings (amplitude anatomy vs. renormalization
subtraction). The synthesis paper should present these as one ``Hopf algebra of
decomposition'' chapter.
\end{corollary}

% =====================================================================
\section{Holographic quantum error correction: strengthening the bridge}
\label{sec:holographic}

The slogan ``gauge redundancy is to physics what fault-tolerant encoding is to quantum
information'' (\Cref{cor:literal-bridge} makes it a theorem for surface codes) is, in
full generality, status \Hh. The most rigorous body of work strengthening it toward
status \Ss\ is \emph{holographic quantum error correction}. We summarize it and label
its epistemic status carefully.

\subsection{Bulk locality as an error-correcting code}

In the AdS/CFT correspondence a bulk (gravitational) operator is redundantly encoded on
the boundary conformal field theory. Almheiri, Dong, and Harlow~\cite{ADH} observed that
this redundancy is precisely that of a quantum error-correcting code: a bulk local
operator in the interior can be reconstructed on any of several boundary subregions,
exactly as a logical operator of a code has representatives on several physical
subsystems~\cite{DongHarlowWall}. Access to
only a boundary subregion is like having part of the code erased; the operator can still
be reconstructed provided the subregion contains the relevant entanglement wedge, a fact
that can be derived from bulk modular flow~\cite{EGH}.

\begin{principle}[Holographic code dictionary, status \Ss/\Hh]
\label{prin:holo}
Under the entanglement-wedge reconstruction of AdS/CFT:
\[
\begin{array}{lcl}
\text{bulk effective field theory} &\leftrightarrow& \text{logical Hilbert space }\Hlog,\\
\text{boundary CFT} &\leftrightarrow& \text{physical Hilbert space }\Hphys,\\
\text{bulk gauge/diffeomorphism redundancy} &\leftrightarrow& \text{code redundancy},\\
\text{radial/holographic direction} &\leftrightarrow& \text{redundancy depth},\\
\text{Ryu--Takayanagi entropy} &\leftrightarrow& \text{code subalgebra / entropy}.
\end{array}
\]
The Pastawski--Yoshida--Harlow--Preskill (HaPPY) tensor-network
codes~\cite{HaPPY} are explicit toy models realizing this dictionary.
\end{principle}

\begin{remark}[Harlow's theorem: RT from QEC]
\label{rem:harlow}
Harlow~\cite{HarlowRT} proved that a quantum-corrected Ryu--Takayanagi formula holds in
\emph{any}
operator-algebra quantum error-correcting code (subalgebra codes,
\Cref{rem:oaqec}), and that when the boundary subalgebras have nontrivial center the
code describes a \emph{gauge theory with edge modes}. This is the sharpest known sense in
which bulk gauge redundancy \emph{is} the redundancy of a code: the center of the
subalgebra is the gauge (superselection/edge-mode) data, exactly the $\Hgauge$ factor of
\eqref{eq:subsystem}. Status: \Ss\ as a theorem about OAQEC codes; \Ss/\Hh\ as a claim
about physical quantum gravity, where finite-$N$ corrections may modify the picture.
\end{remark}

\subsection{Epistemic assessment}

\begin{itemize}
\item For \textbf{surface/topological codes}, ``gauge redundancy $=$ QHA'' is a theorem
(\Cref{cor:literal-bridge}): status \Ss.
\item For \textbf{OAQEC/subalgebra codes with center}, the identification ``gauge (edge
mode) data $=$ subalgebra center'' is a theorem (Harlow): status \Ss.
\item For \textbf{physical AdS/CFT quantum gravity}, the holographic-code picture is
strongly supported but subject to finite-$N$/gauge-invariance caveats: status \Ss/\Hh.
\item The fully general slogan across \emph{all} gauge theories remains a guiding
heuristic: status \Hh.
\end{itemize}
We therefore report the bridge as a spectrum from theorem (surface codes) to heuristic
(fully general), rather than as a single claim, exactly the discipline the series'
status system was designed to enforce.

% =====================================================================
\section{Results}
\label{sec:results}

We collect the principal results and their epistemic status.

\begin{theorem}[Summary of Part VI]
\label{thm:summary}
Let $(\cat C, J)$ be a site, $G$ a group acting on $X$, and consider stabilizer/surface
codes as above. Then:
\begin{enumerate}
\item \textup{(T1, \Cref{thm:aut-stab}, status \Ss)}
$\Aut_{\quot{X}{G}}(x)\cong\Stab_G(x)$, and $\quot{X}{G}$ is strictly more faithful than
$X/G$ exactly on the isotropy $H^\bullet_{\mathrm{grp}}(\Stab_G(x))$.
\item \textup{(T2, \Cref{thm:stackification}, status \Ss; application \Ss/\Hh)} the
representation prestack $\Repp$ admits a universal stackification $\Repp^{J}$, which is
Part I's representation stack; the ladder closes (\Cref{cor:closure}).
\item \textup{(T3, \Cref{thm:surface-code}, status \Ss)} surface-code $Z$- and $X$-type
logical operators are $H_1(\Sigma;\ZZ_2)$ and $H^1(\Sigma;\ZZ_2)$, each $\cong\ZZ_2^{2g}$,
so $N(S)/S\cong\ZZ_2^{4g}$ encodes $2g$ logical qubits; boundaries are pure gauge. Hence
gauge redundancy $=$ QHA is a theorem for surface codes (\Cref{cor:literal-bridge}).
\item \textup{(T4, \Cref{thm:connes-kreimer}, status \Ss)} the Connes--Kreimer antipode
computes the BPHZ counterterm and renormalization is the Birkhoff decomposition of the
Feynman-rules character; Parts II and VI share one Hopf-algebra backbone
(\Cref{cor:shared-hopf}).
\end{enumerate}
\end{theorem}

\begin{proof}
Each part is proved in the referenced statement: (1) \Cref{thm:aut-stab},
\Cref{cor:faithfulness}; (2) \Cref{thm:stackification}, \Cref{cor:closure}; (3)
\Cref{thm:surface-code}, \Cref{cor:literal-bridge}; (4) \Cref{thm:connes-kreimer},
\Cref{cor:shared-hopf}.
\end{proof}

\begin{table}[t]
\centering
\renewcommand{\arraystretch}{1.25}
\begin{tabular}{@{}>{\raggedright\arraybackslash}p{4.4cm}%
                  >{\raggedright\arraybackslash}p{4.4cm}%
                  >{\raggedright\arraybackslash}p{3.3cm}c@{}}
\hline
\textbf{Mathematics} & \textbf{Physical representation} & \textbf{QI interpretation} &
\textbf{Status}\\
\hline
Presheaf & local observable assignment & local syndrome data & \Ss\\
Sheaf & fields from local data & syndrome consistency (local gluing) & \Ss\\
Descent & coordinate-independent physics & correctable global error / recovery & \Ss\\
Stack $\quot{X}{G}$ & gauge field w/ redundancy & encoding $+$ equivalences & \Ss\\
$\Aut_{\quot{X}{G}}(x)\!\cong\!\Stab_G(x)$ & residual/isotropy symmetry & error-detecting
redundancy & \Ss\\
Classifying stack $\mathrm{B}G$ & gauge-field classifier & code-family classifier & \Ss\\
$H_1(\Sigma;\ZZ_2)$ & conserved cycles & surface-code logical qubits & \Ss\\
Boundary $c=\partial b$ & pure gauge / null & trivial (in-stabilizer) error & \Ss\\
Hopf coproduct $\cop$ & decomposition & nested-error bookkeeping & \Ss\\
Antipode $\antipode$ & renormalization inverse & counterterm & \Ss/\Hh\\
Gauge $=$ QHA (general) & redundancy $=$ fault tolerance & holographic code & \Hh\\
\hline
\end{tabular}
\caption{Part VI representation dictionary with status labels (excerpt of Appendix F,
augmented with quantum-information column). Composite labels mark entries standard as
mathematics but heuristic in their fully general physical reading.}
\label{tab:dictionary}
\end{table}

\FloatBarrier

% =====================================================================
\section{Discussion}
\label{sec:discussion}

\subsection{How Part VI composes on Parts I--V}

Part VI is the module with the most cross-references, as befits a capstone:
\begin{itemize}
\item \textbf{From Part V}~\cite{ML-partV} it takes fibered categories, groupoids, and the
$(2,1)$-categorical machinery needed for \Cref{def:stack} and
\Cref{thm:stackification}; univalence (Part V) is the type-theoretic shadow of the
descent equivalence \eqref{eq:descent-equiv}.
\item \textbf{To Part I}~\cite{ML-partI} it delivers \Cref{thm:stackification}, discharging
the forward reference and closing the ladder (\Cref{cor:closure}).
\item \textbf{From Part IV}~\cite{ML-partIV} it reuses homology: \Cref{thm:surface-code}
makes Part IV's dictionary entry literal.
\item \textbf{From Part II}~\cite{ML-partII} it reuses the Hopf-algebraic coproduct:
\Cref{thm:connes-kreimer} is Part II's decomposition axiom applied to divergences
(\Cref{cor:shared-hopf}).
\item \textbf{From Part III}~\cite{ML-partIII} it reuses moduli stacks: $\quot{X}{G}$ is
the gauge-theoretic special case of Part III's moduli stacks.
\end{itemize}

\subsection{Limitations}

We reproduce and address the series' limitations (Part I, \S0.5):
\begin{enumerate}
\item \emph{Not every analogy is a theory.} Status labels are part of the formalism:
\Cref{tab:dictionary} carefully distinguishes the status-\Ss\ surface-code identity
(\Cref{cor:literal-bridge}) from the status-\Hh\ general slogan (\Cref{sec:holographic}).
\item \emph{Gauge redundancy $\neq$ physical symmetry} (\Cref{rem:limitation-gauge}). The
stack $\quot{X}{G}$ is the right object \emph{only} when $G$ is declared pure
description; the same mathematics with $G$ a physical symmetry would give a genuinely
different physics. This modeling choice is not decided by the mathematics.
\item \emph{No single universal functor $\Math\to\Phys$.} \Cref{thm:stackification} is
\emph{local} (indexed by the site $\Dom$); it does not assert one global functor. This
respects the series' deliberate domain-indexed design.
\item \emph{Finite-$N$ caveats.} The holographic-code bridge (\Cref{sec:holographic}) is
subject to gauge-invariance and finite-$N$ corrections; we label it \Ss/\Hh\ accordingly
and do not overclaim it as a theorem about physical quantum gravity.
\end{enumerate}

\subsection{Open problems}

\begin{itemize}
\item Beyond surface codes and OAQEC, is there a general theorem making ``gauge
redundancy $=$ fault-tolerant encoding'' status \Ss\ for an arbitrary gauge theory?
The holographic literature (Harlow; entanglement-wedge reconstruction) is the most
promising route.
\item Can the stackification $\Repp^{J}$ (\Cref{thm:stackification}) be carried out
explicitly for a concrete site of physical theories (factorization homology; Schreiber's
cohesive-topos program)?
\item Is there a single graded connected Hopf algebra unifying the Goncharov and
Connes--Kreimer coproducts (\Cref{cor:shared-hopf})?
\end{itemize}

% =====================================================================
\section{Conclusion}
\label{sec:conclusion}

Part VI has formalized sheaves, stacks, gauge redundancy, and quantum high-availability
as one construction seen four ways: \emph{many local representatives, one invariant
logical content, plus explicit data recording the redundancy}. We proved that the
quotient stack retains isotropy (T1), that the representation prestack stackifies to
Part I's representation stack and thereby closes the modular ladder (T2), that
surface-code logical operators are literally $H_1(\Sigma;\ZZ_2)$ so that ``gauge
redundancy $=$ quantum high-availability'' is a theorem for surface codes (T3), and that
Connes--Kreimer renormalization is Hopf-algebraic decomposition sharing Part II's
coalgebraic backbone (T4). We surveyed holographic quantum error correction as the
rigorous frontier strengthening the general bridge, with each claim's epistemic status
(\Ss/\Hh/\Pp) explicitly labeled.

The four closing theses of the series specialize here to:
\begin{quote}\itshape
symmetry redundancy is extra descriptive freedom that does not change physical
information; a gauge field is a physical field represented by locally redundant
descriptions; a stack is the mathematical object that remembers local descriptions,
gluing, and redundancy symmetries; quantum high availability is logical information
protected by encoding it across redundant physical degrees of freedom.
\end{quote}
Gauge redundancy is to physics what fault-tolerant encoding is to quantum information:
proven for surface codes, conjectured in general, and now formalized as the capstone
module of the representation library. The synthesis paper (Part VII) will assemble the
closure theorem (\Cref{cor:closure}) with the four master formulas and the status-label
statistics into a single recomposition of the ladder.

% =====================================================================
\section*{Companion formalizations}
\addcontentsline{toc}{section}{Companion formalizations}

% ragged right here avoids loose interword spacing around unbreakable code paths
\raggedright
The constructions above are encoded in companion code.

\begin{itemize}
\item \textbf{Haskell} (\path{src/sheaves-stacks-gauge-quantum-ha/}): modules
\texttt{Sheaf}, \texttt{Stack} (coarse vs.\ stack quotient, returning the retained
stabilizer), \texttt{QHA} (encoding, syndrome, recovery), \texttt{Stabilizer}
(surface-code stabilizers and $H_1$ logicals), and \texttt{Hopf} (Connes--Kreimer
coproduct and antipode). The \texttt{Main} module demonstrates T1 (isotropy retained),
T3 (genus-$g$ code dimension $2^{2g}$), and T4, whose antipode passes the Hopf axiom
$m(\antipode\otimes\id)\cop=0$.

\item \textbf{Lean} (\path{lean/sheaves-stacks-gauge-quantum-ha/}): a Mathlib-free,
best-effort sketch of the types \texttt{SheafOnSite}, \texttt{StackOnSite},
\texttt{QuotientStack}, and \texttt{StabilizerCode}, together with the theorem signatures
\texttt{aut\_equiv\_stabilizer} (T1) and \texttt{surfaceCode\_logical\_eq\_h1Rank} (T3);
the file type-checks, with the substantive proofs recorded as deferred \texttt{sorry}
goals in the idiomatic sketch style.
\end{itemize}

% =====================================================================
\begin{thebibliography}{99}
\raggedright

\bibitem{KashiwaraSchapira}
M.~Kashiwara and P.~Schapira,
\emph{Sheaves on Manifolds},
Grundlehren der mathematischen Wissenschaften \textbf{292},
Springer-Verlag, Berlin, 1990 (corrected reprint 1994).

\bibitem{Giraud}
J.~Giraud,
\emph{Cohomologie non ab\'elienne},
Grundlehren der mathematischen Wissenschaften \textbf{179},
Springer, 1971.

\bibitem{Vistoli}
A.~Vistoli,
``Notes on Grothendieck topologies, fibered categories and descent theory,''
arXiv:math/0412512 (2004).

\bibitem{Kitaev}
A.~Yu.~Kitaev,
``Fault-tolerant quantum computation by anyons,''
arXiv:quant-ph/9707021 (1997);
\emph{Ann. Physics} \textbf{303} (2003), 2--30.

\bibitem{DKLP}
E.~Dennis, A.~Kitaev, A.~Landahl, and J.~Preskill,
``Topological quantum memory,''
arXiv:quant-ph/0110143;
\emph{J. Math. Phys.} \textbf{43} (2002), 4452--4505.

\bibitem{Gottesman}
D.~Gottesman,
``Stabilizer Codes and Quantum Error Correction,''
Caltech Ph.D. thesis, arXiv:quant-ph/9705052 (1997).

\bibitem{ConnesKreimer}
A.~Connes and D.~Kreimer,
``Renormalization in quantum field theory and the Riemann--Hilbert problem I:
The Hopf algebra structure of graphs and the main theorem,''
arXiv:hep-th/9912092;
\emph{Comm. Math. Phys.} \textbf{210} (2000), 249--273.

\bibitem{Goncharov}
A.~B.~Goncharov,
``Galois symmetries of fundamental groupoids and noncommutative geometry,''
arXiv:math/0208144;
\emph{Duke Math. J.} \textbf{128} (2005), 209--284.

\bibitem{KnillLaflamme}
E.~Knill and R.~Laflamme,
``Theory of quantum error-correcting codes,''
\emph{Phys. Rev. A} \textbf{55} (1997), 900--911;
arXiv:quant-ph/9604034.

\bibitem{ADH}
A.~Almheiri, X.~Dong, and D.~Harlow,
``Bulk locality and quantum error correction in AdS/CFT,''
arXiv:1411.7041;
\emph{J. High Energy Phys.} \textbf{04} (2015), 163.

\bibitem{HaPPY}
F.~Pastawski, B.~Yoshida, D.~Harlow, and J.~Preskill,
``Holographic quantum error-correcting codes: Toy models for the bulk/boundary
correspondence,''
arXiv:1503.06237;
\emph{J. High Energy Phys.} \textbf{06} (2015), 149.

\bibitem{HarlowRT}
D.~Harlow,
``The Ryu--Takayanagi formula from quantum error correction,''
arXiv:1607.03901;
\emph{Comm. Math. Phys.} \textbf{354} (2017), 865--912.

\bibitem{DongHarlowWall}
X.~Dong, D.~Harlow, and A.~C.~Wall,
``Reconstruction of bulk operators within the entanglement wedge in gauge-gravity
duality,''
arXiv:1601.05416;
\emph{Phys. Rev. Lett.} \textbf{117} (2016), 021601.

\bibitem{EGH}
T.~Faulkner and A.~Lewkowycz (and related),
``Bulk locality from modular flow,''
arXiv:1704.05464;
\emph{J. High Energy Phys.} \textbf{07} (2017), 151.

\bibitem{RotaBaxter}
K.~Ebrahimi-Fard, L.~Guo, and D.~Kreimer,
``Integrable renormalization II: The general case,''
\emph{Ann. Henri Poincar\'e} \textbf{6} (2005), 369--395;
arXiv:hep-th/0403118.

\bibitem{ML-partI}
M.~Long and The YonedaAI Collaboration,
``Foundations: The Representation Stack and Realization Pipeline'' (Part I of
\emph{A \MtoP{} Representation Library}), YonedaAI Research Collective, 2026.

\bibitem{ML-partII}
M.~Long and The YonedaAI Collaboration,
``Motives, Periods, and Amplitudes'' (Part II), YonedaAI Research Collective, 2026.

\bibitem{ML-partIII}
M.~Long and The YonedaAI Collaboration,
``Algebraic Geometry: Spaces of Physical Possibility'' (Part III), YonedaAI Research
Collective, 2026.

\bibitem{ML-partIV}
M.~Long and The YonedaAI Collaboration,
``Algebraic Topology: Conserved Global Information'' (Part IV), YonedaAI Research
Collective, 2026.

\bibitem{ML-partV}
M.~Long and The YonedaAI Collaboration,
``Category Theory and Homotopy Type Theory: Composition and Identity'' (Part V),
YonedaAI Research Collective, 2026.

\end{thebibliography}

\end{document}
